%%%%%%%%%%%%%%%%%%%%%%%%%%%%%%%%%%%%%%%%%%%%%%%%%%%%%%%%%%%%%%%%%%%%%%%
%%%%%%%%%%%%%%%%%%%%%%%%%%%%%%%%%%%%%%%%%%%%%%%%%%%%%%%%%%%%%%%%%%%%%%%
%%%%%%%%%%%%%%%%%%%%%%%%%%%%%%%%%%%%%%%%%%%%%%%%%%%%%%%%%%%%%%%%%%%%%%%

\section{Conditional expectation}

For two random variables $X$ and $Y$, let $\ExCond{X}{Y}$ denote the
expected value of $X$, if the value of $Y$ is specified. Formally, we
have
\begin{equation*}
    \ExCond{X}{Y =y}%
    =%
    \sum_{x \in \Omega} x \ProbCond{X =x}{Y=y}.
\end{equation*}
The expression $\ExCond{X}{Y}$, which is a shorthand for
$\ExCond{X}{Y=y}$, is the \emphi{conditional expectation} of $X$ given
$Y$. In reality, it is a function from the value of $y$, to the
average value of $X$. As such, one can think of conditional
expectation as a function $f(y) = \ExCond{X}{Y=y}$.


\begin{fragment}{conditional:expectation}
\begin{lemma}%
    \lemlab{conditional:expectation}%
    % 
    For any two random variables $X$ and $Y$, we have
    \begin{math}
        \Ex{ \bigl. \ExCond{X}{ Y} } = \Ex{ \bigl. X}.
    \end{math}
\end{lemma}
\end{fragment}
\begin{proof}
    \begin{math}
        \Ex{\bigl. \ExCond{X }{ Y} }%
        =%
        \ExExt{Y}{ \Ex{ X \sep{ Y = y}}} =%
        \sum_y \Prob{Y = y} \ExCond{ X}{ Y = y}%
    \end{math}
    \begin{align*}
      &=%
        \sum_y \Prob{\bigl. Y = y} \frac{\sum_{x} x \Prob{X = x \cap
        Y=y}}{\Prob{Y=y}}%
      \\ &%
           =%
           \sum_{y} \sum_{x} x \Prob{ \bigl. X =x \cap Y = y}%
           =%
           \sum_{x} x \sum_y \Prob{ \bigl. X =x \cap Y = y}%
      \\ &%
           =%
           \sum_{x} x \Prob{\bigl. X = x}%
           =%
           \Ex{\bigl. X}.
    \end{align*}
\end{proof}

\begin{lemma}
    \lemlab{expectations:product}%
    % 
    For any two random variables $X$ and $Y$, we have
    \begin{math}
        \Ex{\bigl.Y \cdot \ExCond{X }{Y}} = \Ex{\bigl. X Y}.
    \end{math}
\end{lemma}
\begin{proof}
    We have that
    \begin{math}
        \ds%
        \Ex{Y \cdot \Ex{X \sep{Y}}}%
        =%
        \sum_y \Prob{Y = y} \cdot y \cdot \Ex{X \sep{Y =y }}
    \end{math}
    \begin{equation*}
        =%
        \sum_y \Prob{Y = y} \cdot y \cdot \frac{\sum_{x} x \Prob{X = x
              \cap Y=y}}{\Prob{Y=y}}%
        =%
        \sum_x \sum_y x y \cdot \Prob{ \bigl. X = x \cap Y = y}%
        =%
        \Ex{\bigl. X Y}.
    \end{equation*}
\end{proof}


\section{\QuickSort{} runs in $O(n \log n)$ time with high
   probability}

Consider a set $T$ of the $n$ items to be sorted, and consider a
specific element $t \in T$. Let $X_i$ be the size of the input in the
$i$\th level of recursion that contains $t$. We know that $X_0 = n$,
and
\begin{align*}
    \Ex{X_i \sep{ X_{i-1}}}%
    \leq%
    \frac{1}{2} \frac{3}{4} X_{i-1} + \frac{1}{2} X_{i-1}%
    \leq%
    \frac{7}{8} X_{i-1}.
\end{align*}
Indeed, with probability $1/2$ the pivot is the middle of the
subproblem; that is, its rank is between $X_{i-1}/4$ and
$(3/4)X_{i-1}$ (and then the subproblem has size $\leq X_{i-1}(3/4)$),
and with probability $1/2$ the subproblem might has not shrank
significantly (i.e., we pretend it did not shrink at all).

Now, observe that for any two random variables we have that
$\Ex{X \MakeSBig } = \ExExt{y}{\Ex{X\left| \MakeSBig {Y=y}
      \right.}\MakeBig }$, see \lemref{conditional:expectation}. As
such, we have that
\begin{align*}
    \Ex{X_i \MakeBig }%
    =%
    \ExExt{y}{\Ex{X_i\sep{X_{i-1}=y}}\Bigl. }%
    \leq%
    \ExExt{X_{i-1}=y}{ \frac{7}{8} y }%
    =%
    \frac{7}{8} \Ex{X_{i-1}\MakeBig }%
    \leq%
    \pth{\frac{7}{8}}^i \Ex{X_0}%
    =%
    \pth{\frac{7}{8}}^i n.
\end{align*}

In particular, consider $M = 8 \log_{8/7} n$. We have that
\begin{align*}
    \mu%
    =%
    \Ex{\MakeBig X_M}%
    \leq%
    \pth{\frac{7}{8}}^M n%
    \leq%
    \frac{1}{n^8} n%
    =%
    \frac{1}{n^7}.
\end{align*}

Of course, $t$ participates in more than $M$ recursive calls, if and
only if $X_M \geq 1$. However, by Markov's inequality
(\thmref{Markov}), we have that
\begin{align*}
    \Prob{ \begin{array}{c}
           \text{element } t \text{ participates }\\
           \text{ in more than } M \text{ recursive calls}
       \end{array}
    }%
    \leq %
    \Prob{ \MakeBig X_M \geq 1 }%
    \leq%
    \frac{\Ex{X_M}}{1} \leq \frac{1}{n^7},
\end{align*}
as desired. That is, we proved that the probability that any element
of the input $T$ participates in more than $M$ recursive calls is at
most $n(1/n^7) \leq 1/n^6$.

\begin{theorem}
    For $n$ elements, \QuickSort runs in $O(n \log n)$ time, with high
    probability.
\end{theorem}

%%% Local Variables:
%%% mode: latex
%%% TeX-master: t
%%% End:
